\documentclass{article}

\usepackage{amsfonts}
\usepackage{amsthm}
\usepackage{mathrsfs}
\usepackage{amsmath}
\usepackage{graphicx}
\newtheorem{theorem}{Theorem}
\theoremstyle{definition}\newtheorem{definition}{DEF}
\theoremstyle{lemma}\newtheorem{lemma}{Lemma}
\theoremstyle{corrolary}\newtheorem{corrolary}{Cor}
\setlength{\parindent}{0pt}

\begin{document}

\begin{definition}
Given a random variable $X$ with cdf $F$, say the expectation is
\begin{equation}
\mathbb{E}(X) = \int_\mathbb{R} x dF
\end{equation}
\end{definition}

\begin{definition}
Given a random variable $X$ with cdf $F$ and a function $g: \mathbb{R} \rightarrow \mathbb{R}$, say
\begin{equation}
\mathbb{E}(g(X)) = \int_\mathbb{R} g dF
\end{equation}
\end{definition}

\textbf{Note-} The expected value does not always exist, it only exists when the (Lebesgue) integral that defines it exists.
For example, let $X$ be a random variable with a Cauchy distribution.
Then it has pdf $f(x)=\frac{1}{1+x^2}$.
This means the expected value is $\int_{-\infty}^\infty \frac{x}{1+x^2} dx$.
This is an integral of an odd function, so we can consider
\begin{align*}
\mathbb{E}(X) &= \int_{-\infty}^\infty \frac{x}{1+x^2} dx \\
&= \int_0^\infty \frac{x}{1+x^2} dx - \int_0^\infty \frac{x}{1+x^2} dx \\
&= \frac{1}{2} \ln (1+x^2) |_0^\infty - \frac{1}{2} \ln (1+x^2) |_0^\infty \\
&= \infty - \infty
\end{align*}
Thus, the expected value is not defined.

\begin{lemma}
The expected value is defined if and only if $\int_\mathbb{R} |x| dF < \infty$
\end{lemma}
\begin{proof}
We say a measurable function is integrable on a measure space when its absolute value has a finite integral.
So the lemma is true by definition.
To build intuition, consider that $x \leq |x|$, so $\int x dF \leq \int |x| dF < \infty$.
\end{proof}

\begin{definition}
Say the mean of a random variable $X$ is $\mu = \mathbb{E}(X)$.
\end{definition}

\textbf{Ex-} Find the mean of a random variable with pdf $f(x) = \lambda e^{-\lambda x}$ on $(0,\infty)$.
Consider 
\begin{align*}
\int_\mathbb{R} |x| dF &= \int_0^\infty x e^{-\lambda x} dx \\
&= xe^{-\lambda x} - \frac{1}{\lambda} e^{-\lambda x} |_0^\infty \\
&= 0 - \frac{-1}{\lambda}\\
&= \frac{1}{\lambda}\\
&< \infty
\end{align*}
Thus, the expected value exists and $\mathbb{E}(X)=\frac{1}{\lambda}$.

\textbf{Ex-} Find the mean of a random variable with pdf $f(k)=\frac{\lambda^k e^{-\lambda}}{k!}$ on $\mathbb{N}$.
\begin{align*}
\int_\mathbb{R} |k| dF &= \sum_{k=0}^\infty \frac{k\lambda^k e^{-\lambda}}{k!} \\
&= \sum_{k=1}^\infty \frac{\lambda^k e^{-\lambda}}{(k-1)!}\\
&= \lambda \sum_{k=1}^\infty \frac{k\lambda^{k-1} e^{-\lambda}}{(k-1)!!}\\
&= \lambda\\
&<\infty
\end{align*}
Thus, the expected value exists and is $\lambda$.

\begin{theorem}
Let $X$ and $Y$ be random variables and $g: \mathbb{R} \rightarrow \mathbb{R}$.
If $\mathbb{E}(Y)$ exists then $\mathbb{E}(Y)=\mathbb{E}(g(X))$.
\end{theorem}
\begin{proof}
By definition of expectation,
\begin{align*}
\mathbb{E}(Y) &= \int_\mathbb{R} y dF_y\\
&= \int_\mathbb{R} g d(F \circ g)\\
&= \mathbb{E}(g(X))
\end{align*}
\end{proof}
\begin{corrolary}
$\mathbb{E}(a+bX) = a + b\mathbb{E}(X)$
\end{corrolary}

\begin{definition}
Say the variance of a random variable with mean $\mu$ is
\begin{equation}
Var(X)=\mathbb{E}((X-\mu)^2)
\end{equation}
\end{definition}

\begin{theorem}
\textbf{Parallell Axis Theorem:}
Let $X$ be a random variable with mean $\mu$ and $c\in \mathbb{R}$.
\begin{equation}
\mathbb{E}((X-c)^2) = Var(X) + (\mu - c)^2
\end{equation}
\end{theorem}
\begin{proof}
\begin{align*}
\mathbb{E}((X-c)^2) &= \int_\mathbb{R} (x-c)^2 dF\\
&= \int_\mathbb{R} (x-\mu +\mu -c)^2 dF\\
&= \int_\mathbb{R} (x-\mu )^2 + 2(x-\mu )(\mu -c) + (\mu -c)^2 dF\\
&= \int_\mathbb{R} (x-\mu )^2 dF + 2 (\mu -c) \int_\mathbb{R} (x-\mu ) dF + \int_\mathbb{R}(\mu -c)^2 dF\\
&= Var(X) + 0 +(\mu -c)^2
\end{align*}
\end{proof}
\begin{corrolary}
$Var(X) = \mathbb{E}(X^2) - \mathbb{E}(X)^2$
\end{corrolary}

\begin{theorem}
\begin{equation}
Var(a + bX) = b^2 Var(X)
\end{equation}
\end{theorem}
\begin{proof}
Let $\mu = \mathbb{E}(a + bX)$.
Then $\mathbb{E}(a + bX) = a + b\mathbb{E}(X)$.
So
\begin{align*}
Var(a+bX) &= \int_\mathbb{R} (a + bx - a - b\mu )^2 dF\\
&= \int_\mathbb{R} (bx - b\mu )^2 dF \\
&= b^2 \int_\mathbb{R} (x-\mu )^2 dF\\
&= b^2 Var(X)
\end{align*}
\end{proof}

\begin{definition}
Say the standard deviation of a random variable is
\begin{equation}
\sigma = \sqrt{Var(X)}
\end{equation}
\end{definition}

\textbf{Ex-} Find the variance of a random variable with odf $f(x)=\lambda e^{-\lambda x}$.
First, remember that $\mathbb{E}(X) = \frac{1}{\lambda}$.
Also,
\begin{align*}
\mathbb{E}(X^2) &= \int_\mathbb{R} (x-\mu)^2 dF\\
&= \int_0^\infty x^2 \lambda e^{-\lambda x}dx\\
&= x^2 e^{-\lambda x} + \frac{2x}{\lambda} e^{-\lambda x} + \frac{2}{\lambda^2} |_0^\infty \\
&= 0 - \frac{-2}{\lambda^2}\\
&= \frac{2}{\lambda^2}
\end{align*}
Thus, $Var(X) = \mathbb{E}(X^2) - \mathbb{E}(X)^2 = \frac{2}{\lambda^2} - \frac{1}{\lambda^2} = \frac{1}{\lambda^2}$.

\begin{definition}
Say the skewness of a random variable $X$ with mean $\mu$ and standard deviation $\sigma$ is
\begin{equation}
skew(X)=\frac{\mathbb{E}((X-\mu)^3)}{\sigma^3}
\end{equation}
\end{definition}

\begin{lemma}
\begin{equation}
skew(X)=\frac{\mathbb{E}(X^3)-3\mu \sigma^2 -\mu^3}{\sigma^3}
\end{equation}
\end{lemma}
\begin{proof}
We only need consider the numerator of the definition.
\begin{align*}
\mathbb{E}((X-\mu)^3) &= \mathbb{E}(X^3) -3\mu \mathbb{E}(X^2) + 3\mu^2 \mathbb{E}(X) -\mu^3 \\
&= \mathbb{E}(X^3) -3\mu (\mu^2 + \sigma^2) + 3\mu^3 -\mu^3\\
&= \mathbb{E}(X^3) -3\mu \sigma^2 -\mu^3
\end{align*}
\end{proof}

\textbf{Ex-} If $X \sim \mathcal{N}(\mu ,\sigma^2 )$, then $skew(X)=0$.

\begin{definition}
Say the kurtosis of a random variable $X$ with mean $\mu$ and standard deviation $\sigma$ is
\begin{equation}
kurt(X) = \frac{\mathbb{E}((X-\mu)^4)}{\sigma^4}
\end{equation}
\end{definition}

\begin{lemma}
The kurtosis is
\begin{equation}
kurt(X) = \frac{\mathbb{E}(X^4)-4\mu \mathbb{E}(X^3)+6\mu^2 \sigma^2 +3\mu^4}{\sigma^4}
\end{equation}
\end{lemma}
\begin{proof}
We only need to consider the numerator.
\begin{align*}
\mathbb{E}((X-\mu)^4) &= \mathbb{E}(X^4) -4\mu \mathbb{E}(X^3) + 6\mu^2 \mathbb{E}(X^2) - 4\mu^3 \mathbb{E}(X)+\mu^4\\
&= \mathbb{E}(X^4) -4\mu \mathbb{E}(X^3) + 6\mu^2 (\mu^2 + \sigma^2) -4\mu^4 + \mu^4\\
&= \mathbb{E}(X^4) -4\mu \mathbb{E}(X^3) + 6\mu^2 \sigma^2 + 3\mu^4
\end{align*}
\end{proof}

\textbf{Ex-} If $X\sim \mathcal{N}(\mu ,\sigma^2)$ then $kurt(X)=3$.

\begin{definition}
Say the excess kurtosis is
\begin{equation}
kurt_{ex}(X)=kurt(X)-3
\end{equation}
\end{definition}

\begin{theorem}
Let $X$ be a random variable and $n\in \mathbb{N}$.
If $\mathbb{E}(X^n)$ exists then $\mathbb{E}(X^k)$ exists for all $k\leq n$.
\end{theorem}
\begin{proof}
Notice that if $|x|<1$, then $|x^k|<1$.
Similarly, if $|x|>1$, then $|x^k|<|x^n|$.
Thus,
\begin{align*}
\int_\mathbb{R} |x^k|dF &= 2 \int_0^\infty |x^k|dF\\
&= 2\int_0^1 |x^k| dF + 2\int_1^\infty |x^k| dF \\
&\leq 2\int_0^1 1dF + 2\int_1^\infty |x^n|dF\\
&\leq \int_\mathbb{R} 1dF + \int_\mathbb{R} |x^n|dF\\ &= 1 + \mathbb{E}(X^2)\\
&<\infty
\end{align*}
\end{proof}

\begin{definition}
Say $\mathbb{E}(X^n)$ is the $n^{th}$ moment of $X$.
\end{definition}

\begin{definition}
Let $X$ be a random variable and $T>0$ such that $\mathbb{E}(e^{Xt})$ exists for $|t|< T$.
Say the moment generating function (MGF) is
\begin{equation}
M_X(t)=\mathbb{E}(e^{Xt})
\end{equation}
\end{definition}

\begin{theorem}
Let $X$ be a random variable with MGF $M_X$.
\begin{equation}
M_X(t)=\sum_{n=0}^\infty \frac{\mathbb{E}(X^n)t^n}{n!}
\end{equation}
\end{theorem}
\begin{proof}
This is an application of Taylor's Theorem, so we need to show $M_X^{(n)}(t)=\mathbb{E}(X^ne^{tX})$.
We will do so via induction.
Notice that when $n=0$, this is the definition of hte mgf.
So the base case is trivial. \\
Now, suppose $M_X^{(k)}(t) = \mathbb{E}(X^ke^{tX})$ for some $k\in \mathbb{N}$.
Then
\begin{align*}
M_X^{(k+1)}(t) &= \frac{d}{dt}\mathbb{E}(X^k e^{tX})\\
&= \mathbb{E}(\frac{\partial}{\partial t}X^ke^{tx})\\
&= \mathbb{E}(X^{k+1}e^{tX})
\end{align*}
Therefore, we have found the derivative via induction.
Finally, we can compute $M_X^{(n)}(0)=\mathbb{E}(X^ne^0)=\mathbb{E}(X^n)$.
Therefore, by Taylor's theorem, we have the result.
\end{proof}

\begin{theorem}
Let $X$ and $Y$ be two random variables with MGFs $M_X$ and $M_Y$.
$M_X = M_Y$ if and only if $X$ and $Y$ have the same distibution.
\end{theorem}
\begin{proof}
First, suppose $X$ and $Y$ have the same distibution.
Then clearly $\mathbb{E}(e^{Xt})=\mathbb{E}(e^{Yt})$.
So $M_X=M_Y$.\\
Conversely, suppose $X$ and $Y$ do not have the same distribution.
So $f_X \neq f_Y$ on a set of non-zero measure.
Akso, $M_X(t)=\int_\mathbb{R} e^{tx}f_x(x)dx$ and $M_Y(t)=\int_\mathbb{R} e^{tx}f_Y(x)dx$.
Then $M_X(t)-M_Y(t)=\int_\mathbb{R} e^{tx}(f_X(x)-f_Y(x))dx$.
Thus, by the zero integral lemma, this integral is non-zero.
Therefore, the MGFs must be distinct.
\end{proof}

\textbf{Ex-} Find the MGF of a random variable with PDF $\lambda e^{-\lambda x}$.
\begin{align*}
M_X(t)&=\mathbb{E}(e^{tX})\\
&=\int_0^\infty e^{tx} \lambda e^{-\lambda x} dx\\
&= \int_0^\infty \lambda e^{(t-\lambda)x}dx\\
&= \frac{\lambda}{t-\lambda} e^{(t-\lambda)x}|_0^\infty\\
&= 0-\frac{\lambda}{t-\lambda}\\
&= \frac{\lambda}{\lambda -t}
\end{align*}
Notice the evaluation of the integral relies on $e^{(t-\lambda)x} \rightarrow 0$ as $x \rightarrow \infty$.
This is true if and only if $t-\lambda < 0$, or $t < \lambda$.
Thus, the MGF is $\frac{\lambda}{\lambda-t}$ with domain $t<\lambda$.

\begin{theorem}
Let $X$ be a random variable with MGF $M(t)$ for $|t|<T$.
Then the MGF of $a+bX$ is $e^{ta}M(bt)$ for $|t|<\frac{T}{|b|}$.
\end{theorem}
\begin{proof}
\begin{align*}
M_{a+bX}(t) &= \mathbb{E}(e^{t(a+bX)})\\
&= e^{ta}\mathbb{E}(e^{tbX})\\
&= e^{ta} M(tb)
\end{align*}
Also, it is clar that$t<\frac{T}{|b|}$ if and only if $|tb|<T$.
\end{proof}

\begin{definition}
Let $X$ be a random variable with mean $\mu$ and let $\xi>0$.
Say $X$ is a $\xi$-subgaussian random variable when
\begin{equation}
\mathbb{E}(e^{t(X-\mu)})\leq e^\frac{t^2\sigma^2}{2}
\end{equation}
for all $t\in \mathbb{R}$.
Say $X$ is subgaussian if it is $\xi$-subgaussian for some $\xi>0$.
\end{definition}

\end{document}
